%% kapa: exam=1 round=22 year=2011 period=2 form=A subject=appraisal_law subject_ko=부동산관계법규 seq=15 number=55
%% source: Q-Net 감정평가사 기출문제 (cst00309) artlSeq=5036790 file=extracted/1008/5036790/2011년 제22회 감정평가사 1차시험 기출문제/2교시.hwp
%% conversion: kordoc 4.24.51 -> tools/gen1.py (auto); review=proofread-gpt-astra; flags=none
\begin{question}{55}
국토의 계획 및 이용에 관한 법령상 이행강제금과 관련하여 옳지 않은 것은?
\begin{choices}
  \item 토지거래계약허가를 받은 자가 허가받은 목적대로 토지를 이용하지 아니할 경우 이용의무의 이행명령을 담보하는 수단이 될 수 있다.
  \item 토지거래계약허가를 받아 휴게음식점을 취득하여 실제로 이용하는 자가 해당 건축물의 일부를 임대하는 경우에는 이행강제금의 부과대상이 아니다.
  \item 이행명령을 받은 자가 그 명령을 이행하더라도 명령을 이행하기 전에 이미 부과된 이행강제금은 징수하여야 한다.
  \item 부과권자는 시장ㆍ군수 또는 구청장이며 이행명령이 정하여진 기간 내에 이행되지 아니할 때 부과하게 된다.
  \item 최초의 이행명령이 있었던 날을 기준으로 하여 6개월에 한 번씩 그 이행명령이 이행될 때까지 반복하여 부과ㆍ징수할 수 있다.
\end{choices}
\end{question}
\answer{5}
